\subsection{Extension of derivations: the case of fields}

Let K, L and M be fields such that \(K \subset L \subset M\) . By Prop. 21 (III, p.~572) we have an exact sequence of vector M-spaces

\[
(E _ {K, L, M})
\]

\[
\Omega_ {\mathrm{K}} (\mathrm{L}) \otimes_ {\mathrm{L}} \mathbf {M} \xrightarrow {\alpha} \Omega_ {\mathrm{K}} (\mathbf {M}) \xrightarrow {\beta} \Omega_ {\mathrm{L}} (\mathbf {M}) \to 0;
\]

the M-linear mappings \(\alpha\) and \(\beta\) are characterized by the relations

\begin{enumerate}
\def\labelenumi{(\arabic{enumi})}
\setcounter{enumi}{2}
\tightlist
\item
\end{enumerate}

\[
\begin{array}{r l} \alpha (d _ {\mathrm{L/K}} x \otimes 1) = d _ {\mathrm{M/K}} x & \text { for } \quad x \in \mathrm{L} \\ \beta (d _ {\mathrm{M/K}} y) = d _ {\mathrm{M/L}} y & \text { for } \quad y \in \mathrm{M}. \end{array}\tag{4}
\]

Let V be a vector M-space, denote by \(D_{K}(M,V)\) the vector space of K-derivations of M with values in V and similarly introduce \(D_{K}(L,V)\) and \(D_{L}(M,V)\) . By III, p.~571, Diagram (42) and III, p.~572, Diagram (44), we have a commutative diagram of homomorphisms of vector spaces

\[
\begin{array}{l}0 \rightarrow \operatorname{Hom} _ {\mathbf {M}} (\Omega_ {\mathrm{L}} (\mathbf {M}), \mathrm{V}) \xrightarrow {\operatorname{Hom} (\beta , 1)} \operatorname{Hom} _ {\mathbf {M}} (\Omega_ {\mathrm{K}} (\mathbf {M}), \mathrm{V}) \xrightarrow {\operatorname{Hom} (\alpha , 1)} \operatorname{Hom} _ {\mathbf {M}} (\Omega_ {\mathrm{K}} (\mathrm{L}) \otimes_ {\mathrm{L}} \mathrm{M}, \mathrm{V})\\\downarrow\\0 \rightarrow \operatorname{Der} _ {\mathrm{L}} (\mathrm{M}, \mathrm{V}) \xrightarrow {i _ {\mathrm{V}}} \operatorname{Der} _ {\mathrm{K}} (\mathrm{M}, \mathrm{V}) \xrightarrow {r _ {\mathrm{V}}} \operatorname{Der} _ {\mathrm{K}} (\mathrm{L}, \mathrm{V}),\end{array}
\]

where the vertical arrows are isomorphisms, \(i_{v}\) is the canonical injection and \(r_{v}\) the restriction mapping.

We thus obtain from II, p.~299, Th. 5 and II, p.~301, Prop. 10 the following proposition:

PROPOSITION 2. --- The following conditions are equivalent :

\begin{enumerate}
\def\labelenumi{\alph{enumi})}
\item
  The mapping \(\alpha\) is injective.
\item
  Every K-derivation of L into M extends to a K-derivation of M into M.
\item
  Every K-derivation of L into a vector M-space V extends to a K-derivation of M into V.
\end{enumerate}

PROPOSITION 3. --- Let K be a field, L a pure extension of K and \((x_{i})_{i\in I}\) a pure basis of L (V, p.~106, Def. 2).

\begin{enumerate}
\def\labelenumi{\alph{enumi})}
\item
  Let \(\mathbf{V}\) be a vector space over \(\mathbf{L}\) , \(\Delta\) a derivation of \(\mathbf{K}\) into \(\mathbf{V}\) and \((v_i)_{i\in I}\) a family of elements of \(\mathbf{V}\) . There exists a unique derivation \(\mathbf{D}\) of \(\mathbf{L}\) into \(\mathbf{V}\) , extending \(\Delta\) and such that \(\mathbf{D}(x_i) = v_i\) for all \(i\in I\) .
\item
  The vector L-space \(\Omega_{\mathbf{K}}(\mathbf{L})\) of K-differentials of L has the family \((dx_i)_{i\in I}\) as basis.
\end{enumerate}

Assertion b) has been proved in IV, p.~23 and Assertion a) follows directly from this.

COROLLARY. --- Let P be a subfield of K. The canonical mapping \(\alpha\) of \(\Omega_{\mathrm{P}}(\mathrm{K}) \otimes_{\mathrm{K}} \mathrm{L}\) into \(\Omega_{\mathrm{P}}(\mathrm{L})\) is injective and the family \((d_{\mathrm{L}/\mathrm{P}}x_{i})_{i \in \mathrm{I}}\) is a basis (over L) of a subspace of \(\Omega_{\mathrm{P}}(\mathrm{L})\) supplementary to the image of \(\alpha: \Omega_{\mathrm{P}}(\mathrm{K}) \otimes_{\mathrm{K}} \mathrm{L} \to \Omega_{\mathrm{P}}(\mathrm{L})\) .

Prop. 3, a) shows that every P-derivation of K into a vector space V over L extends to a P-derivation of L into V; the injectivity of \(\alpha\) then follows by Prop. 2. The second assertion of the corollary follows from the exactness of the sequence \((E_{P,K,L})\) and Prop. 3, b) (taking Formula (4) into account).

PROPOSITION 4. --- Let K be a field, L a separable algebraic extension of K and V a vector space over L.

\begin{enumerate}
\def\labelenumi{\alph{enumi})}
\item
  Every K-derivation of L into V is zero.
\item
  If \(\Delta\) is a derivation of \(K\) into \(V\) , there exists a unique derivation \(D\) of \(L\) into \(V\) which extends \(\Delta\) .
\end{enumerate}

Let D be a K-derivation of L into V. If E is a subextension of L of finite degree over K, then the K-algebra E is etale, and so \(\Omega_{\mathrm{K}}(\mathrm{E}) = 0\) (V, p.~33, Th. 3), whence \(D \mid E = 0\) by the universal property of \(\Omega_{\mathrm{K}}(\mathrm{E})\) (III, p.~569). Since L is the union of a family of subextensions of finite degree over K, we have D = 0, whence a).

Let \(\Delta\) be a derivation of K into V. If D' and D'' are two extensions of \(\Delta\) to a derivation of L into V, then the difference D' - D'' is a K-derivation of L into V; hence it is zero by \(a\) ) and so D' = D''.

It remains to prove the existence of an extension of \(\Delta\) . Zorn's Lemma (E, III, p.~20) implies the existence of a maximal extension \(D_0\) of \(\Delta\) to a derivation defined on a subfield \(L_0\) of \(L\) containing \(K\) .

Let \(x\) be in \(L\) and \(g\) the minimal polynomial of \(x\) over \(L_0\) . Since \(L\) is algebraic and separable over \(K\) , \(x\) is algebraic and separable over \(L_0\) (V, p.~39, Prop. 6 and p.~39, Cor. 2); therefore \(x\) is a simple root of \(g\) (V, p.~39, Prop. 5), whence \(g'(x) \neq 0\) (IV, p.~17, Prop. 7). If we define \(g^{D_0}(x)\) as in \(V\) , p.~127, there exists thus an element \(u\) of \(V\) such that \(g^{D_0}(x) + g'(x) \cdot u = 0\) ; by Prop. 1 (V, p.~127) there exists a derivation \(D\) of \(L_0(x)\) into \(V\) extending \(D_0\) and such that \(D(x) = u\) . In view of the maximal character of \((L_0, D_0)\) we thus have \(L_0(x) = L_0\) , whence \(x \in L_0\) . Since \(x\) was arbitrary we conclude that \(L_0 = L\) .

COROLLARY 1. --- We have \(\Omega_{\mathbf{K}}(\mathbf{L}) = 0\) if \(\mathbf{L}\) is algebraic and separable over \(\mathbf{K}\) .

The corollary follows from Prop. 4, a) because the vector L-space \(\Omega_{\mathbf{K}}(\mathbf{L})\) is generated by the image of the canonical K-derivation \(d_{\mathrm{L / K}}: \mathrm{L} \to \Omega_{\mathbf{K}}(\mathrm{L})\) .

COROLLARY 2. --- If L is an algebraic and separable extension of a field K, the canonical mapping \(\alpha: \Omega_{\mathrm{P}}(\mathrm{K}) \otimes_{\mathrm{K}} \mathrm{L} \to \Omega_{\mathrm{P}}(\mathrm{L})\) is an isomorphism for every subfield P of K.

The mapping \(\alpha\) is injective by Prop. 2 (V, p.~128) and Prop. 4, \(b\) ); since \(\Omega_{\mathrm{K}}(\mathrm{L})\) reduces to 0 (Cor. 1) the exactness the sequence \((\mathrm{E}_{\mathrm{P},\mathrm{K},\mathrm{L}})\) implies that \(\alpha\) is surjective.

COROLLARY 3. --- Let E be an extension of a field K and D a derivation of E into E, mapping K into K. If L is a subextension of E which is algebraic and separable over K, then D(L) ⊂ L.

Let \(\Delta\) be the derivation of K into L which agrees with D on K. By Prop. 4 (V, p.~129) there exists a derivation \(D'\) of L into L extending \(\Delta\) . Now we may consider \(D'\) and the restriction \(D''\) of D to L as derivations of L into E; since they agree on K, we have \(D' = D''\) by Prop. 4, whence

\[
\mathrm{D} (\mathrm{L}) = \mathrm{D} ^ {\prime \prime} (\mathrm{L}) = \mathrm{D} ^ {\prime} (\mathrm{L}) \subset \mathrm{L}.
\]

Remarks. --- 1) Later (V, p.~131, Cor. 3 and p.~135, Cor. 2) we shall prove a converse to Cor. 1 of Prop. 4.

\begin{enumerate}
\def\labelenumi{\arabic{enumi})}
\setcounter{enumi}{1}
\tightlist
\item
  Every algebraic extension of a prime field is separable (V, p.~37, Cor.). Since every derivation of a prime field is zero, every derivation of an algebraic extension of a prime field is zero (Prop. 4).
\end{enumerate}

